\section{Развој ARCore апликације}

Пре имплементације самог АР прегледа учионица потребно је припремити развојно окружење, обезбедити уређај који подржава ARCore и поставити структуру пројекта са неопходним зависностима. У наставку су описани ти кораци, при чему наведене верзије алата и библиотека одговарају стварном пројекту из овог рада.

\subsection{Потребна опрема и софтвер на рачунару}

Развој се одвија у окружењу \textit{Android Studio}, званичном интегрисаном развојном окружењу за Android платформу. Довољно је инсталирати последњу стабилну верзију, јер она већ садржи одговарајући JDK (верзије 17 или новије) и менаџер за \textit{Android SDK}, кроз који се преузима платформа која одговара вредности \texttt{compileSdk = 36} из овог пројекта. Програмски језик је \textit{Kotlin} 2.1.20, у комбинацији са \textit{Android Gradle Plugin}-ом 8.13.0, док се сам \textit{Gradle} не инсталира посебно --- пројекат садржи \textit{Gradle wrapper}, скрипту која при првом покретању преузима тачно ону верзију алата са којом је пројекат тестиран.

Када је реч о емулатору, развој АР апликације реално захтева физички уређај: емулатор подржава ARCore само у ограниченим сценаријима, са виртуелном сценом уместо стварног окружења, па се понашање апликације --- посебно препознавање одштампаног маркера правом камером и стабилност сидрења --- не може веродостојно проверити без правог телефона.

\subsection{Потребна опрема на телефону}

За покретање апликације неопходан је Android уређај сертификован за ARCore --- Google одржава јавну листу подржаних модела чије су камере, сензори покрета и процесори прошли проверу квалитета праћења. Уређај мора имати Android 7.0 (API ниво 24) или новији, што се поклапа са вредношћу \texttt{minSdk = 24} у конфигурацији пројекта. Поред тога, на уређају мора бити инсталирана услуга \textit{Google Play Services for AR} (ARCore), која се дистрибуира и аутоматски ажурира преко продавнице \textit{Google Play}, па апликација не пакује ARCore у себе, већ се ослања на ту системску услугу.

Да би се апликација инсталирала директно из окружења \textit{Android Studio}, на телефону је потребно укључити опције за програмере и дозволити отклањање грешака преко USB везе (\textit{USB debugging}); алтернативно, новије верзије окружења подржавају и бежично упаривање преко Wi-Fi мреже. Пошто апликација податке преузима са сервера у локалној мрежи, телефон и рачунар на коме сервер ради морају бити повезани на исту Wi-Fi мрежу.

\subsection{Креирање пројекта и зависности}

Пројекат се креира кроз чаробњак у окружењу \textit{Android Studio}, избором шаблона \textit{Empty Views Activity}, са језиком \textit{Kotlin} и конфигурацијом изградње у облику \textit{Gradle Kotlin DSL} датотека (\texttt{build.gradle.kts}). Након креирања, у датотеку \texttt{app/build.gradle.kts} додају се зависности приказане у табели~\ref{tab:dependencies}, чије су верзије централизоване у каталогу \texttt{gradle/libs.versions.toml}.

\begin{table}[H]
\centering
\small
\begin{tabular}{lr}
\toprule
\textbf{Зависност} & \textbf{Верзија} \\
\midrule
\texttt{io.github.sceneview:arsceneview} & 2.2.1 \\
\texttt{androidx.core:core-ktx} & 1.16.0 \\
\texttt{androidx.appcompat:appcompat} & 1.7.0 \\
\texttt{com.google.android.material:material} & 1.12.0 \\
\texttt{androidx.constraintlayout:constraintlayout} & 2.2.1 \\
\texttt{androidx.lifecycle:lifecycle-runtime-ktx} & 2.8.7 \\
\texttt{com.squareup.retrofit2:retrofit} & 2.11.0 \\
\texttt{com.squareup.retrofit2:converter-kotlinx-serialization} & 2.11.0 \\
\texttt{org.jetbrains.kotlinx:kotlinx-serialization-json} & 1.8.0 \\
\texttt{com.squareup.okhttp3:logging-interceptor} & 4.12.0 \\
\bottomrule
\end{tabular}
\caption{Зависности коришћене у пројекту}
\label{tab:dependencies}
\end{table}

Кључна зависност је библиотека \textit{SceneView} (\texttt{arsceneview}), која обједињује ARCore и \textit{Filament} механизам за рендеровање: уместо директног коришћења ARCore-а, што би захтевало ручно писање OpenGL кода за приказ камере и 3Д објеката, SceneView нуди готову компоненту \texttt{ARSceneView}, граф сцене са чворовима (\texttt{Node}) и рендеровање „из кутије". За мрежну комуникацију са сервером користи се \textit{Retrofit} у комбинацији са \textit{kotlinx-serialization} конвертором, чиме се JSON одговори сервера аутоматски претварају у типизиране Kotlin објекте, док се \textit{OkHttp logging-interceptor} укључује само у развојној варијанти ради увида у мрежни саобраћај. Корутине и \texttt{lifecycle-runtime-ktx} обезбеђују асинхроно периодично преузимање података усклађено са животним циклусом активности. У блоку \texttt{android} укључена је и опција \texttt{buildFeatures \{ viewBinding = true \}}, којом се за сваки XML распоред генерише типизована класа за приступ елементима интерфејса, чиме се избегавају позиви \texttt{findViewById}.

\subsection{Подешавање манифеста}

Датотека \texttt{AndroidManifest.xml} мора прецизно описати захтеве апликације према систему. Пројекат декларише следеће ставке:

\begin{itemize}
\item Дозвола за приступ камери, која је обавезна, јер без камере проширена стварност није могућа, и дозвола за приступ интернету, неопходна за преузимање података са сервера:

\begin{verbatim}
<uses-permission android:name="android.permission.CAMERA" />
<uses-permission android:name="android.permission.INTERNET" />
\end{verbatim}

\item Декларација да је АР хардвер обавезан, у комбинацији са ознаком да је ARCore неизоставна компонента апликације:

\begin{verbatim}
<uses-feature
    android:name="android.hardware.camera.ar"
    android:required="true" />
<meta-data
    android:name="com.google.ar.core"
    android:value="required"
    tools:replace="android:value" />
\end{verbatim}

Овим паром апликација постаје „AR Required", што значи да се у продавници нуди само уређајима који подржавају ARCore. Ознака \texttt{tools:replace} над атрибутом \texttt{android:value} овде је неопходна: библиотека SceneView у свом манифесту исту мета-ознаку декларише са вредношћу \texttt{optional}, па би спајање манифеста (\textit{manifest merge}) без ове ознаке завршило грешком због сукоба вредности; овако вредност из апликације надјачава вредност из библиотеке.

\item Атрибут \texttt{android:usesCleartextTraffic="true"} на елементу \texttt{<application>}, којим се дозвољава нешифрован HTTP саобраћај. Од Android-а 9 систем подразумевано блокира саобраћај без TLS-а, а сервер из овог рада ради у локалној мрежи на адреси облика \texttt{http://192.168.0.49:8080}, без сертификата, па је изузетак неопходан за демонстрацију.
\end{itemize}

\subsection{Кораци до прве АР сцене}

Када су окружење, зависности и манифест спремни, до прве функционалне АР сцене долази се кроз неколико корака:

\begin{enumerate}
\item У XML распоред активности додаје се компонента \texttt{ARSceneView}, која истовремено приказује слику са камере и рендероване 3Д објекте.
\item Поглед на сцену повезује се са животним циклусом активности доделом \texttt{sceneView.lifecycle = lifecycle}, чиме се АР сесија аутоматски паузира и наставља заједно са активношћу.
\item У позиву \texttt{configureSession} конфигурише се сесија: гради се база слика \texttt{AugmentedImageDatabase} са маркерима учионица и додељује пољу \texttt{config.augmentedImageDatabase}, чиме се укључује праћење слика описано у одељку 2.5.
\item Подразумевани приказ детектованих равни искључује се поставком својства \texttt{isEnabled} компоненте \texttt{planeRenderer} на вредност нетачно, јер апликацији равни нису потребне, а истакнуте површине би визуелно оптерећивале екран.
\item Преко повратног позива \texttt{onSessionUpdated} апликација се укључује у обраду сваког фрејма и из њега чита новопрепознате и ажуриране слике маркера.
\item У време извршавања тражи се дозвола за камеру, будући да се од API нивоа 23 опасне дозволе морају експлицитно одобрити од стране корисника.
\item Апликација се покреће на физичком уређају, са одштампаним маркером, где се проверава да ли се маркер препознаје и да ли панел стабилно стоји изнад њега.
\end{enumerate}

У наредном поглављу приказано је како се од ових основних градивних елемената --- сесије, базе слика, сидара и чворова --- гради комплетна функционалност АР прегледа учионица.
